\scanpage{75}
\mspage{J}
% archive foliation: 74
car $>0$. Cela prouve $r=0$, et achève la démonstration.

\uline{Corollaire 1} Invariance birationnelle de $\uline{\mathrm{Pic}}^\tau_{X/k}$
pour $X$ propre et simple sur $k$.

\uline{Corollaire 2} Soit $K$ une extension
\deleted{régulière}\inserted{\unsure{séparable primaire|séparablement primaire}}
de type fini de $k$, admettant un modèle \deleted{régulier}
propre et \uline{simple} sur $k$. Alors
$\varinjlim_X\uline{\mathrm{Pic}}^\tau_{X/k}$, où $X$ parcourt
les modèles propres et normaux \inserted{``sur $k$''} de $K$
existe au sens usuel [le système \deleted{projectif}\inserted{inductif}
est essentiellement constant].

\margin{\uline{N.B.} Si $k$ n'est pas parfait, et $K$ est de dim $1$ sur $k$,
\deleted{\unsure{donc}} c'est une condition un peu idiote!}

\uline{Question}. Le cor. 2 reste-t-il valable
en supposant seulement que $K$ admette
un modèle ``normal sur $k$''??
\deleted{\unsure{Il y aurait}}\inserted{\unsure{Cela}}
\deleted{\unsure{du moins}} équivaudrait à la conjonction des
assertions suivantes
\[
\begin{array}{ll}
\text{(i)}&\dim\uline{\mathrm{Pic}}_{X/k}\text{ reste majoré}\\
\text{(ii)}&\text{ordre }\boxed{\mathrm{Tors}\,\uline{\mathrm{Pic}}(X/k)}
                    \text{ reste majoré}\\
\text{(iii)}&\dim\boxed{H^1(X,\underline O_X)}\text{ reste majoré.}
\end{array}
\]
\margin{N.B. (i) résulte de (iii)}
